% !TEX root = ../../presentation.tex
% =============================================================
%  TikZ — the line as a two-port: terminals and the coordinate x
%  Terminal quantities U_S/I_S and U_R/I_R, and the local quantities
%  U/I at distance x along a line of total length l.
%  The pi-equivalent of the same line is fig/pi-section.tex.
%  Voltage arrows are inset so they do not touch the terminals.
%  Self-contained snippet -- \input directly inside a frame.
%  Scales to fill \linewidth, so the enclosing column (or a
%  minipage) sets the size.
% =============================================================
\resizebox{\linewidth}{!}{%
\begin{tikzpicture}[
    x=0.60cm, y=0.60cm,
    line width=0.9pt,
    every node/.style={inner sep=1.5pt},
    term/.style={circle, draw, fill=white, line width=0.9pt,
                 minimum size=4.5pt, inner sep=0pt},
    flow/.style={-{Triangle[length=2.8mm,width=2.4mm]}, line width=0.9pt},
    dim/.style={dotted, line width=1pt},
  ]

  % ---- geometry ---------------------------------------------
  \def\xa{0}     % sending-end terminals
  \def\xb{8}     % receiving-end terminals
  \def\xm{2.9}   % cross-section at distance x
  \def\yt{4}     % upper conductor
  \def\yb{0}     % lower conductor

  % ---- conductors -------------------------------------------
  \draw (\xa,\yt) -- (\xb,\yt);
  \draw (\xa,\yb) -- (\xb,\yb);

  % ---- line current I at the cross-section ------------------
  \draw[flow] (\xa,\yt) -- (\xm,\yt);
  \node[above=1pt] at (\xm,\yt) {$\underline{I}$};

  % ---- terminals --------------------------------------------
  \node[term] at (\xa,\yt) {};
  \node[term] at (\xb,\yt) {};
  \node[term] at (\xa,\yb) {};
  \node[term] at (\xb,\yb) {};

  % ---- terminal currents ------------------------------------
  \draw[flow] ({\xa-1.7},\yt) -- ({\xa-0.8},\yt);
  \node[left] at ({\xa-1.9},\yt) {$\underline{I}_S$};
  \draw[flow] ({\xb+0.8},\yt) -- ({\xb+1.7},\yt);
  \node[right] at ({\xb+1.9},\yt) {$\underline{I}_R$};

  % ---- voltages (upper conductor -> lower conductor) --------
  \draw[flow] (\xa,{\yt-0.4}) -- (\xa,{\yb+0.4});
  \node[right=2pt] at (\xa,{(\yt+\yb)/2}) {$\underline{U}_S$};
  \draw[flow] (\xm,{\yt-0.4}) -- (\xm,{\yb+0.4});
  \node[right=2pt] at (\xm,{(\yt+\yb)/2}) {$\underline{U}$};
  \draw[flow] (\xb,{\yt-0.4}) -- (\xb,{\yb+0.4});
  \node[right=2pt] at (\xb,{(\yt+\yb)/2}) {$\underline{U}_R$};

  % ---- coordinate x and overall length l --------------------
  \node[below=1pt] at (\xm,\yb) {$x$};
  \draw[dim] (\xa,{\yb-0.4}) -- (\xa,{\yb-2.4});
  \draw[dim] (\xb,{\yb-0.4}) -- (\xb,{\yb-2.4});
  \draw[dim] (\xa,{\yb-1.9}) -- (\xb,{\yb-1.9});
  \node[below=1pt] at ({(\xa+\xb)/2},{\yb-2.0}) {$l$};

\end{tikzpicture}%
}
